\chapter{Discussion}


\todo{Iets over scrum}


At the start of this project, we were introduced to the scrum method of working together. In the beginning, we tried using this method, as we divided the defined scrum roles. We quickly realized that the division of the tasks weren't as effective as they could have been in larger projects as the tasks only consisted of programming assignments and some theory segments. To do the theory segments before the programming tasks did not make a lot of sense and we decided to only do the theory after the programming tasks. Nonetheless, as there is a chronological order to the assignments, we ended up doing one programming assignment at a time. So eventually we did not continue the method, as we didn't adhere to the Scrum framework by organizing scrum sprints. The reason for this was that we estimated that it would not be effective in such a small and short project. We came to our own method, wherein we conducted occasional meetings. In these meetings, we discussed everyone's progress and determined what was to be done next. An overall planning was also maintained. This worked quite well, however, in hindsight, we should have conducted more meetings to work together more effectively. Now 2 situations will be described in which Scrum could have benefited.

The first situation occurred halfway through the project. The midterm deadline was coming up, and the group was split in two teams of two people. While writing the report, this complicated things, as we weren't well aware what the other team was doing. This led to the two parts of the midterm report being written in a completely different style. This was not good, as the report therefore was less coherent and harder to read. Next time, it would be useful to incorporate sprint reviews of the scrum method, allowing us to be better up to date with what everyone is doing, and therefore coherently work together. 

The second situation occurred multiple times during our project. While working on the project, a lot of code was developed. In most cases, this code also had to be adjusted in one way or another, which was doable as we implemented an efficient code structure. However, not all team members were aware of this structure and therefore didn't adhere to this. This resulted in very different coding styles, which led to more time being spent in comprehending what everyone else had made. This could be solved by a more dedicated product owner (from the scrum method), who manages the code bases and makes sure everyone understands the standards that are being adhered to.








